
\begin{tikzpicture}[
    scale=0.7,
    transform shape,
    node distance=7mm,
    font=\small,
    >=Stealth,
    % Dinh dang chung
    base/.style={
        draw=cyan!70!blue,
        fill=cyan!12,
        thick,
        align=center,
        inner sep=5pt
    },
    % Bat dau / Ket thuc: hinh elip
    startstop/.style={
        base,
        ellipse,
        minimum width=3.3cm,
        minimum height=0.85cm
    },
    % Nhap du lieu: hinh binh hanh
    input/.style={
        base,
        trapezium,
        trapezium left angle=70,
        trapezium right angle=110,
        text width=9.4cm,
        minimum height=1.05cm
    },
    % Xu ly / tinh toan: hinh chu nhat
    block/.style={
        base,
        rectangle,
        text width=10.2cm,
        minimum height=1.05cm
    },
    % Dieu kien: hinh thoi
    decision/.style={
        base,
        diamond,
        aspect=2.3,
        text width=2.8cm,
        inner sep=1.5pt
    },
    % Xuat du lieu: hinh tai lieu
    output/.style={
        base,
        tape,
        tape bend top=none,
        text width=10.2cm,
        minimum height=1.3cm
    },
    % Mui ten luong xu ly
    arr/.style={
        ->,
        thick,
        draw=cyan!75!blue
    },
    % Nhan cua nhanh dieu kien
    branch/.style={
        font=\small,
        fill=white,
        inner sep=2pt
    }
]

% ==========================================
% 1. BAT DAU
% ==========================================

\node[startstop] (start) {
    \textbf{Bắt đầu}
};

% ==========================================
% 2. NHAN DU LIEU
% ==========================================

\node[input, below=of start] (data) {
    \textbf{Nhận dữ liệu GY-85 có dấu thời gian}\\
    Gia tốc, vận tốc góc và từ trường
};

% ==========================================
% 3. HIEU CHINH DU LIEU
% ==========================================

\node[block, below=of data] (calib) {
    \textbf{Hiệu chỉnh và chuẩn hóa dữ liệu}\\
    Hiệu chỉnh bias, từ trường;\\
    đổi từ hệ trục cảm biến sang hệ trục lắp đặt
};

% ==========================================
% 4. XU LY GIA TOC
% ==========================================

\node[block, below=of calib] (accel) {
    \textbf{Xử lý tín hiệu gia tốc}\\
    $a=\sqrt{a_x^2+a_y^2+a_z^2}$\\
    Lọc thành phần phù hợp;
    kiểm tra đỉnh và khoảng thời gian tối thiểu
};

% ==========================================
% 5. UOC LUONG TU THE VA HUONG
% ==========================================

\node[block, below=of accel] (heading) {
    \textbf{Ước lượng tư thế và hướng}\\
    Hợp nhất gyro, gia tốc và từ trường hợp lệ;\\
    đổi hướng sang hệ mặt bằng,
    chuẩn hóa góc và trạng thái tin cậy
};

% ==========================================
% 6. KIEM TRA BUOC DI
% ==========================================

\node[decision, below=10mm of heading] (stepcheck) {
    \textbf{Phát hiện}\\
    \textbf{bước hợp lệ?}
};

% ==========================================
% 7. UOC LUONG CHIEU DAI BUOC
% ==========================================

\node[block, below=10mm of stepcheck] (stride) {
    \textbf{Ước lượng chiều dài bước $\ell_k$}\\
    Chỉ sử dụng chiều dài bước mới phát hiện,\\
    không sử dụng tổng quãng đường tích lũy
};

% ==========================================
% 8. CAP NHAT TOA DO PDR
% ==========================================

\node[block, below=of stride] (pdr) {
    \textbf{Cập nhật vị trí tại sự kiện bước $k$}\\[2pt]
    $x_k=x_{k-1}+\ell_k\sin\psi_k$\\[2pt]
    $y_k=y_{k-1}+\ell_k\cos\psi_k$\\[3pt]
    $x$: hướng Đông,\quad $y$: hướng Bắc\\
    $\psi=0$: hướng Bắc
};

% ==========================================
% 9. XUAT DU LIEU
% ==========================================

\node[output, below=of pdr] (send) {
    \textbf{Đóng gói và đồng bộ tới gateway}\\
    Tọa độ, hướng, số bước, tuổi dữ liệu
    và trạng thái chất lượng\\
    Vị trí PDR chỉ là giá trị ước lượng
};

% ==========================================
% 10. KIEM TRA TIEP TUC
% ==========================================

\node[decision, below=10mm of send] (running) {
    \textbf{Tiếp tục}\\
    \textbf{giám sát?}
};

% ==========================================
% 11. KET THUC
% ==========================================

\node[startstop, below=10mm of running] (end) {
    \textbf{Kết thúc}
};

% ==========================================
% NOI CAC LUONG XU LY CHINH
% ==========================================

\draw[arr] (start.south) -- (data.north);
\draw[arr] (data.south) -- (calib.north);
\draw[arr] (calib.south) -- (accel.north);
\draw[arr] (accel.south) -- (heading.north);
\draw[arr] (heading.south) -- (stepcheck.north);

% Co buoc hop le
\draw[arr]
    (stepcheck.south)
    -- node[branch, right] {Có}
    (stride.north);

% Khong phat hien buoc:
% doc mau cam bien tiep theo
\draw[arr]
    (stepcheck.west)
    -- node[branch, above] {Không}
    ++(-4.4cm,0)
    |- (data.west);

% Cap nhat va xuat du lieu
\draw[arr] (stride.south) -- (pdr.north);
\draw[arr] (pdr.south) -- (send.north);
\draw[arr] (send.south) -- (running.north);

% Tiep tuc giam sat:
% quay lai nhan du lieu moi
\draw[arr]
    (running.east)
    -- node[branch, above] {Có}
    ++(4.4cm,0)
    |- (data.east);

% Dung giam sat
\draw[arr]
    (running.south)
    -- node[branch, right] {Không}
    (end.north);

\end{tikzpicture}
